\section*{ \begin{center}
RESUMEN
\end{center} }

\vspace{0.4cm}

\noindent
El entorno universitario contemporáneo obliga a los estudiantes a procesar un volumen exponencial de información técnica no estructurada, lo que desencadena una severa sobrecarga cognitiva e isla los esfuerzos de aprendizaje. Para mitigar estas deficiencias, la presente investigación evalúa el impacto empírico del ecosistema digital SynapTome en el rendimiento académico, la retención y los flujos de trabajo colaborativo de los estudiantes. Metodológicamente, la plataforma integra un pipeline adaptativo de reconocimiento de texto manuscrito (HTR) basado en la arquitectura Transformer TrOCR sometida a un proceso de \textit{fine-tuning} específico, un middleware de Generación Aumentada por Recuperación (RAG) condicionado y acoplado a Modelos de Lenguaje Grande (LLMs), y una topología de Grafo de Conocimiento indexada de forma relacional y vectorial, sincronizada en tiempo real mediante SignalR.

\vspace{0.3cm}

\noindent
Los resultados de campo demostraron un cambio de paradigma drástico para los estudiantes. El pipeline optimizado de TrOCR redujo la tasa de error de caracteres (CER) en la notación científica compleja desde un 34.20\% inicial (empleando herramientas OCR tradicionales) hasta un 4.15\% estable, eliminando la pérdida de información técnica. Asimismo, el middleware RAG restringió el espacio de inferencia semántica al corpus académico validado, mitigando las alucinaciones de los LLMs hasta un nivel asintótico del 1.20\%. Este alineamiento factual elevó las calificaciones promedio en exámenes parciales de un rango regular (11.80) a un nivel sobresaliente (15.90). Finalmente, la generación automatizada de cuestionarios basados en repetición espaciada, en combinación con las interconexiones topológicas, optimizó los procesos cognitivos, reduciendo el estrés académico subjetivo en 46.30 puntos en la escala NASA-TLX y asegurando una retención del conocimiento a largo plazo superior al 93.50\%. La validación estadística mediante una prueba t de Student unilateral derecha ($t = 3.21 > t_{\text{crítico}} = 2.92$) rechazó la hipótesis nula con un nivel de confianza del 95\%, concluyendo que la implementación del ecosistema inteligente SynapTome optimiza significativamente la gestión del aprendizaje, la digitalización de precisión y la ejecución de la inteligencia colectiva entre los estudiantes universitarios en comparación con los métodos aislados tradicionales.

\vspace{0.6cm}

\noindent
\textbf{Palabras clave:} Ecosistema de Gestión del Conocimiento, Generación Aumentada por Recuperación (RAG), Transformer TrOCR, Grafos de Conocimiento, Inteligencia Colectiva, Mitigación de Sobrecarga Cognitiva.